Prostate cancer is one of the most common malignant neoplasms in men, representing a significant public health challenge \cite{ref_siegel2024}. Statistics reported by the International Agency for Research on Cancer indicate that in 2022 this neoplasm was the fourth most common worldwide, accounting for 7.3\% of all cancer cases \cite{JournalStatistic2024}. In this context, tumor grading plays a central role in determining clinical management and prognosis. The Gleason system, based on the sum of the two predominant histological patterns, evolved into the grading system proposed by the International Society of Urological Pathology (ISUP), organized into five groups ranging from 1 to 5. This system provides improved prognostic stratification and is widely adopted in clinical practice \cite{ref-epstein2016}.

The traditional histopathological evaluation of Whole-Slide Images (WSIs) relies on manual analysis performed by pathologists. Given the large number of fields that must be examined within a single slide, the evaluation process becomes a repetitive task that requires long working hours, sustained attention, and specialized expertise. Fatigue resulting from prolonged analysis may introduce inter- and intra-observer variability in results, directly influencing patient clinical management \cite{ref-bulten2020}. Consequently, in recent years, there has been significant growth in the use of Artificial Intelligence (AI) techniques to automate the analysis of histopathological examinations and assist in prostate cancer diagnosis, aiming to mitigate the aforementioned limitations \cite{expert2022review}.

In this scenario, Convolutional Neural Networks (CNNs) have represented the state of the art in medical image analysis \cite{ref-campanella2019}, being capable of extracting relevant features and combining complex visual patterns to identify specific pathological conditions \cite{araujo2021active}. This capability makes the approach promising for the automated analysis of Gleason patterns and ISUP grade groups, thereby reducing the variability inherent to conventional evaluation. Based on this premise, deep learning techniques for prostate grading have been widely investigated in the literature, with notable advances in WSIs, patch extraction strategies, and multi-scale aggregation mechanisms. For example, in the context of weakly supervised learning, \cite{Xiang2023Prostate} proposed the GCN-MIL framework, which uses a Graph Neural Network to aggregate spatial features of patches on WSI slides. The model incorporates a Robust Training strategy that iteratively filters high-uncertainty samples to mitigate the impact of noisy labels, achieving a quadratic kappa of 0.931 on the internal PANDA challenge dataset. However, the approach employs the traditional \textit{Cross-Entropy} loss function, which treats ISUP grades as independent nominal classes and ignores the biological hierarchy of cancer progression.

In an independent clinical validation with 593 images from the \textit{Seoul National University Hospital}, \cite{jung2022artificial} evaluated the DeepDx Prostate system, achieving a quadratic kappa of 0.922 in Gleason grading. However, the model's generalization was limited by the homogeneity of the institutional dataset and by a tendency to overestimate Grade Group 1 (GG1) cases. Seeking indirect risk prediction, \cite{liu2022deep} investigated the use of deep learning to identify patients at risk of cancer from slides previously classified as benign. The method achieved an AUC of 0.739 for risk prediction, revealing subtle morphological patterns, but was limited by reliance on a single institutional domain.

To address the inherent scale variability of glandular structures, \cite{pyramidalcnn} proposed a \textit{Pyramidal CNN} architecture composed of three shallow networks operating on different patch sizes. Although the method achieved approximately 0.77 accuracy in pattern classification, it was limited by semi-automatic labeling and moderate performance in more complex classes. In parallel, aiming to improve interpretability, \cite{kosoko2024xai_prostate} evaluated architectures combined with \textit{Explainable AI} techniques (\textit{Grad-CAM}) using the SICAPv2 dataset. The VGG19 model achieved the best performance (AUC-ROC of 0.85), although the study lacked external clinical validation. Following a simplification strategy, \cite{afifi2024prostate} formulated the detection problem as a binary classification task (benign vs. malignant tissues). Using \textit{fine-tuning} on InceptionResNetV2 models with the \textit{binary cross-entropy} loss, they achieved 91.8\% accuracy, emphasizing detection over multiclass Gleason grading.

In advanced architectures, the incorporation of attention mechanisms has also gained prominence. \cite{alici2024efficient} proposed a deep learning model based on the EfficientNet-B4 architecture with a channel attention mechanism (Eff4-Attn), achieving 93.32\% accuracy in an eight-class classification task using images from the \textit{DiagSet}. However, the study relies exclusively on isolated patches and suffers from class imbalance. Complementarily, \cite{ikromjanov2021vit_prostate} investigated the use of \textit{Vision Transformers} (ViT) on a subset of the PANDA dataset, achieving patch classification accuracy above 85\%, highlighting the ability of attention mechanisms to capture global contextual dependencies.

Despite recent advances, several methodological gaps remain. Many studies treat the problem as a nominal classification task and employ traditional loss functions, disregarding the ordinal nature of ISUP grades. Consequently, prediction errors are penalized uniformly, without accounting for the magnitude of the discrepancy between classes, which does not properly reflect their clinical relevance. Moreover, many approaches do not incorporate robust mechanisms for noise filtering or uncertainty modeling, an important aspect given the presence of artifacts and inconsistencies in public datasets.

To address these limitations, this work proposes a pipeline for the automatic classification of ISUP grade groups in histopathological images. The main contributions are: (i) a hybrid loss function that combines ordinal penalization and focal loss, considering the hierarchical nature of the problem and class imbalance; (ii) a noise removal strategy based on uncertainty estimated from the average predictions of EfficientNet models (B0, B1, B2, and B3); and (iii) a comparative evaluation of EfficientNet architectures B0, B3, and B7, aiming to obtain a more robust and consistent inference model.